Part I. Mathematical Foundations (≈60 pages)
Historical Development
Curie principle
Landau's theory
Superconductivity
Nambu's work
Goldstone theorem
Higgs mechanism
Symmetry in Physics
Groups
Lie groups
Lie algebras
Representations
Noether theorem
Vacuum manifolds
Classical Spontaneous Symmetry Breaking
Double-well potential
Mexican hat potential
Catastrophe theory
Stability analysis
Bifurcations

Part II. Quantum Field Theory (≈120 pages)
Canonical Quantization
Path Integral Formalism
Effective Action
Goldstone Theorem
Including complete proofs
equal-time commutators
Ward identities
current algebra
LSZ derivation
spectral representation
Weinberg proof
Coleman theorem
Mermin–Wagner theorem
Every proof carried out line-by-line.
Higgs Mechanism
Detailed derivation
Abelian Higgs model
Non-Abelian Higgs model
Rξ gauges
BRST symmetry
Fadeev–Popov ghosts
Physical spectrum
Gauge independence
Renormalization
including
one-loop effective potential
Coleman-Weinberg mechanism
beta functions
dimensional regularization
spontaneous symmetry breaking under RG flow

Part III. Condensed Matter (≈80 pages)
Ferromagnetism
Superconductivity
Superfluidity
Bose-Einstein Condensation
Crystals
Liquid crystals
Magnons
Phonons
Anderson-Higgs mechanism
including microscopic derivations from many-body theory.

Part IV. Particle Physics (≈100 pages)
Chiral Symmetry Breaking
NJL Model
QCD Vacuum
Chiral Perturbation Theory
Electroweak Symmetry Breaking
Higgs Sector
Vacuum Stability
Composite Higgs
Technicolor
Little Higgs
Twin Higgs
Dilaton models
Part V. Cosmology (≈80 pages)
Inflation
Phase Transitions
Bubble nucleation
False vacuum decay
Coleman-De Luccia instantons
Topological defects
Cosmic strings
Monopoles
Domain walls
Baryogenesis
Gravitational waves from phase transitions

Part VI. Lorentz Symmetry Breaking (≈80 pages)
Since you have expressed interest in this topic previously, this part could be developed in particular depth.
Chapters would include
Preferred Frames
Einstein-aether theory
Bumblebee models
Ghost condensates
Hořava gravity
Standard Model Extension
Spontaneous Lorentz breaking in string theory
Goldstone modes for spacetime symmetries
Relation to dark matter
Cosmological preferred frames
CMB frame
Global time
Experimental constraints

Part VII. Advanced Topics (≈80 pages)
Supersymmetry breaking
Conformal symmetry breaking
Scale symmetry
Time crystals
Non-equilibrium SSB
Quantum information viewpoint
Tensor network description
Holography
Quantum gravity
Emergent spacetime
Appendices
A. Group Theory
B. Lie Algebras
C. Differential Geometry
D. Functional Methods
E. Path Integrals
F. Distribution Theory
G. Renormalization
H. Topology
I. Homotopy Groups
J. Instantons
K. BRST Quantization
L. Effective Field Theory
M. Large-N Expansion
N. Lattice Field Theory

Proposed Contents
Part I Mathematical Foundations
Chapter 1
Historical Development of Symmetry Breaking
Chapter 2
Mathematical Theory of Symmetry
Chapter 3
Continuous Groups and Lie Algebras
Chapter 4
Noether's Theorem
Chapter 5
Classical Spontaneous Symmetry Breaking

Part II Quantum Field Theory
Chapter 6
Canonical Quantization
Chapter 7
Path Integral Formulation
Chapter 8
Vacuum Structure
Chapter 9
Effective Action
Chapter 10
Goldstone Theorem
Chapter 11
Coleman Theorem
Chapter 12
Higgs Mechanism
Chapter 13
Gauge Symmetry Breaking
Chapter 14
BRST Quantization
Chapter 15
Renormalization
Chapter 16
Coleman–Weinberg Mechanism

Part III Condensed Matter
Chapter 17
Landau Theory
Chapter 18
Ferromagnetism
Chapter 19
Superconductivity
Chapter 20
Superfluidity
Chapter 21
BEC
Chapter 22
Anderson Mechanism

Part IV Particle Physics
Chapter 23
Nambu–Jona-Lasinio Model
Chapter 24
Chiral Symmetry Breaking
Chapter 25
QCD Vacuum
Chapter 26
Chiral Perturbation Theory
Chapter 27
Electroweak Symmetry Breaking
Chapter 28
Composite Higgs

Part V Cosmology
Chapter 29
Finite Temperature Field Theory
Chapter 30
Phase Transitions
Chapter 31
False Vacuum Decay
Chapter 32
Topological Defects
Chapter 33
Inflation
Chapter 34
Baryogenesis
Part VI Gravity
Chapter 35
Lorentz Symmetry Breaking
Chapter 36
Einstein–Aether Theory
Chapter 37
Bumblebee Models
Chapter 38
Preferred Frames
Chapter 39
Quantum Gravity
Chapter 40
Open Problems
with complete derivations.
